% Options for packages loaded elsewhere
% Options for packages loaded elsewhere
\PassOptionsToPackage{unicode}{hyperref}
\PassOptionsToPackage{hyphens}{url}
\PassOptionsToPackage{dvipsnames,svgnames,x11names}{xcolor}
%
\documentclass[
  spanish,
  letterpaper,
  DIV=11,
  numbers=noendperiod]{scrartcl}
\usepackage{xcolor}
\usepackage{amsmath,amssymb}
\setcounter{secnumdepth}{5}
\usepackage{iftex}
\ifPDFTeX
  \usepackage[T1]{fontenc}
  \usepackage[utf8]{inputenc}
  \usepackage{textcomp} % provide euro and other symbols
\else % if luatex or xetex
  \usepackage{unicode-math} % this also loads fontspec
  \defaultfontfeatures{Scale=MatchLowercase}
  \defaultfontfeatures[\rmfamily]{Ligatures=TeX,Scale=1}
\fi
\usepackage{lmodern}
\ifPDFTeX\else
  % xetex/luatex font selection
\fi
% Use upquote if available, for straight quotes in verbatim environments
\IfFileExists{upquote.sty}{\usepackage{upquote}}{}
\IfFileExists{microtype.sty}{% use microtype if available
  \usepackage[]{microtype}
  \UseMicrotypeSet[protrusion]{basicmath} % disable protrusion for tt fonts
}{}
\makeatletter
\@ifundefined{KOMAClassName}{% if non-KOMA class
  \IfFileExists{parskip.sty}{%
    \usepackage{parskip}
  }{% else
    \setlength{\parindent}{0pt}
    \setlength{\parskip}{6pt plus 2pt minus 1pt}}
}{% if KOMA class
  \KOMAoptions{parskip=half}}
\makeatother
% Make \paragraph and \subparagraph free-standing
\makeatletter
\ifx\paragraph\undefined\else
  \let\oldparagraph\paragraph
  \renewcommand{\paragraph}{
    \@ifstar
      \xxxParagraphStar
      \xxxParagraphNoStar
  }
  \newcommand{\xxxParagraphStar}[1]{\oldparagraph*{#1}\mbox{}}
  \newcommand{\xxxParagraphNoStar}[1]{\oldparagraph{#1}\mbox{}}
\fi
\ifx\subparagraph\undefined\else
  \let\oldsubparagraph\subparagraph
  \renewcommand{\subparagraph}{
    \@ifstar
      \xxxSubParagraphStar
      \xxxSubParagraphNoStar
  }
  \newcommand{\xxxSubParagraphStar}[1]{\oldsubparagraph*{#1}\mbox{}}
  \newcommand{\xxxSubParagraphNoStar}[1]{\oldsubparagraph{#1}\mbox{}}
\fi
\makeatother


\usepackage{longtable,booktabs,array}
\usepackage{calc} % for calculating minipage widths
% Correct order of tables after \paragraph or \subparagraph
\usepackage{etoolbox}
\makeatletter
\patchcmd\longtable{\par}{\if@noskipsec\mbox{}\fi\par}{}{}
\makeatother
% Allow footnotes in longtable head/foot
\IfFileExists{footnotehyper.sty}{\usepackage{footnotehyper}}{\usepackage{footnote}}
\makesavenoteenv{longtable}
\usepackage{graphicx}
\makeatletter
\newsavebox\pandoc@box
\newcommand*\pandocbounded[1]{% scales image to fit in text height/width
  \sbox\pandoc@box{#1}%
  \Gscale@div\@tempa{\textheight}{\dimexpr\ht\pandoc@box+\dp\pandoc@box\relax}%
  \Gscale@div\@tempb{\linewidth}{\wd\pandoc@box}%
  \ifdim\@tempb\p@<\@tempa\p@\let\@tempa\@tempb\fi% select the smaller of both
  \ifdim\@tempa\p@<\p@\scalebox{\@tempa}{\usebox\pandoc@box}%
  \else\usebox{\pandoc@box}%
  \fi%
}
% Set default figure placement to htbp
\def\fps@figure{htbp}
\makeatother



\ifLuaTeX
\usepackage[bidi=basic]{babel}
\else
\usepackage[bidi=default]{babel}
\fi
% get rid of language-specific shorthands (see #6817):
\let\LanguageShortHands\languageshorthands
\def\languageshorthands#1{}


\setlength{\emergencystretch}{3em} % prevent overfull lines

\providecommand{\tightlist}{%
  \setlength{\itemsep}{0pt}\setlength{\parskip}{0pt}}



 


\KOMAoption{captions}{tableheading}
\makeatletter
\@ifpackageloaded{caption}{}{\usepackage{caption}}
\AtBeginDocument{%
\ifdefined\contentsname
  \renewcommand*\contentsname{Tabla de contenidos}
\else
  \newcommand\contentsname{Tabla de contenidos}
\fi
\ifdefined\listfigurename
  \renewcommand*\listfigurename{Listado de Figuras}
\else
  \newcommand\listfigurename{Listado de Figuras}
\fi
\ifdefined\listtablename
  \renewcommand*\listtablename{Listado de Tablas}
\else
  \newcommand\listtablename{Listado de Tablas}
\fi
\ifdefined\figurename
  \renewcommand*\figurename{Figura}
\else
  \newcommand\figurename{Figura}
\fi
\ifdefined\tablename
  \renewcommand*\tablename{Tabla}
\else
  \newcommand\tablename{Tabla}
\fi
}
\@ifpackageloaded{float}{}{\usepackage{float}}
\floatstyle{ruled}
\@ifundefined{c@chapter}{\newfloat{codelisting}{h}{lop}}{\newfloat{codelisting}{h}{lop}[chapter]}
\floatname{codelisting}{Listado}
\newcommand*\listoflistings{\listof{codelisting}{Listado de Listados}}
\makeatother
\makeatletter
\makeatother
\makeatletter
\@ifpackageloaded{caption}{}{\usepackage{caption}}
\@ifpackageloaded{subcaption}{}{\usepackage{subcaption}}
\makeatother
\usepackage{bookmark}
\IfFileExists{xurl.sty}{\usepackage{xurl}}{} % add URL line breaks if available
\urlstyle{same}
\hypersetup{
  pdftitle={Inteligencia Artificial, aprendizaje y evaluación del aprendizaje},
  pdfauthor={CAICAS - Analytics},
  pdflang={es},
  colorlinks=true,
  linkcolor={blue},
  filecolor={Maroon},
  citecolor={Blue},
  urlcolor={Blue},
  pdfcreator={LaTeX via pandoc}}


\title{Inteligencia Artificial, aprendizaje y evaluación del
aprendizaje}
\usepackage{etoolbox}
\makeatletter
\providecommand{\subtitle}[1]{% add subtitle to \maketitle
  \apptocmd{\@title}{\par {\large #1 \par}}{}{}
}
\makeatother
\subtitle{Reporte académico ampliado a partir de \texttt{reporte.md}}
\author{CAICAS - Analytics}
\date{2026-06-23}
\begin{document}
\maketitle

\renewcommand*\contentsname{Tabla de contenidos}
{
\hypersetup{linkcolor=}
\setcounter{tocdepth}{3}
\tableofcontents
}

\section{Propósito y alcance}\label{propuxf3sito-y-alcance}

Este reporte complementa y expande el documento fuente
\texttt{reporte.md}, cuyo eje argumental sostiene que la inteligencia
artificial generativa obliga a reconfigurar la evaluación del
aprendizaje, desplazar el foco desde el producto final hacia el proceso
cognitivo, y formar competencias de juicio crítico, agencia humana,
ética y trazabilidad. El presente documento amplía esa tesis con
literatura académica digital sobre IA en educación, aprendizaje asistido
por IA, retroalimentación, evaluación formativa, integridad académica y
rediseño de instrumentos evaluativos.

El criterio de inclusión usado fue restrictivo: se priorizaron fuentes
académicas con DOI verificable, pertinencia explícita con aprendizaje o
evaluación, y disponibilidad digital mediante página editorial,
repositorio institucional, SSRN, Springer, Elsevier, Wiley, Taylor \&
Francis, MDPI, AJET o revistas universitarias. No se incorporaron
documentos institucionales sin DOI dentro de la matriz de referencias,
aunque algunos de ellos puedan ser relevantes para política educativa.

\section{Punto de partida del documento
fuente}\label{punto-de-partida-del-documento-fuente}

El documento base plantea siete ideas principales que se conservan y se
desarrollan en este reporte:

\begin{enumerate}
\def\labelenumi{\arabic{enumi}.}
\tightlist
\item
  La IA generativa no debe entenderse solo como una herramienta
  tecnológica, sino como un factor de reconfiguración curricular,
  pedagógica y evaluativa.
\item
  La evaluación tradicional centrada en textos finales, ensayos
  estáticos o productos aislados pierde validez cuando dichos productos
  pueden ser generados, refinados o simulados por IA.
\item
  La evaluación debe desplazarse hacia evidencias de proceso: bitácoras,
  defensa oral, explicación de decisiones, trazabilidad de fuentes,
  iteraciones y transferencia contextual.
\item
  El uso de IA puede producir un ``rendimiento aparente'': mejora del
  desempeño mientras la herramienta está disponible, sin garantía de
  aprendizaje autónomo posterior.
\item
  La respuesta institucional no debería reducirse a vigilancia o
  detección automática, porque esta estrategia es técnicamente limitada
  y pedagógicamente insuficiente.
\item
  Las competencias estudiantiles y docentes deben incluir alfabetización
  crítica en IA, comprensión de sesgos, privacidad, explicabilidad,
  ética y límites de la automatización.
\item
  La evaluación resiliente a la IA debe certificar autonomía, juicio
  disciplinar y capacidad de transferencia, no solo producción textual.
\end{enumerate}

\section{Metodología de búsqueda y
validación}\label{metodologuxeda-de-buxfasqueda-y-validaciuxf3n}

La búsqueda se orientó por cuatro familias de literatura:

\begin{itemize}
\tightlist
\item
  \textbf{IA en educación superior y aprendizaje:} revisiones
  sistemáticas y metarrevisiones sobre aplicaciones de IA en educación.
\item
  \textbf{IA generativa y aprendizaje:} estudios sobre tutoría, agencia
  estudiantil, retroalimentación, metacognición y riesgos de
  dependencia.
\item
  \textbf{Evaluación del aprendizaje:} literatura fundacional sobre
  evaluación formativa, retroalimentación y autorregulación.
\item
  \textbf{Evaluación en la era de IA generativa:} rediseño de tareas,
  autenticidad, integridad académica, escalas de uso permitido de IA y
  limitaciones de detectores.
\end{itemize}

La validación DOI se realizó verificando que cada identificador tuviera
una página digital resoluble mediante el DOI o estuviera confirmado por
la página editorial, repositorio institucional o registro académico de
la publicación. La tabla siguiente resume la verificación.

\section{Matriz de fuentes académicas digitales con DOI
validado}\label{matriz-de-fuentes-acaduxe9micas-digitales-con-doi-validado}

\begin{longtable}[]{@{}
  >{\raggedleft\arraybackslash}p{(\linewidth - 10\tabcolsep) * \real{0.2000}}
  >{\raggedright\arraybackslash}p{(\linewidth - 10\tabcolsep) * \real{0.1500}}
  >{\raggedleft\arraybackslash}p{(\linewidth - 10\tabcolsep) * \real{0.2000}}
  >{\raggedright\arraybackslash}p{(\linewidth - 10\tabcolsep) * \real{0.1500}}
  >{\raggedright\arraybackslash}p{(\linewidth - 10\tabcolsep) * \real{0.1500}}
  >{\raggedright\arraybackslash}p{(\linewidth - 10\tabcolsep) * \real{0.1500}}@{}}
\toprule\noalign{}
\begin{minipage}[b]{\linewidth}\raggedleft
Nº
\end{minipage} & \begin{minipage}[b]{\linewidth}\raggedright
Fuente
\end{minipage} & \begin{minipage}[b]{\linewidth}\raggedleft
Año
\end{minipage} & \begin{minipage}[b]{\linewidth}\raggedright
Tema principal
\end{minipage} & \begin{minipage}[b]{\linewidth}\raggedright
DOI validado
\end{minipage} & \begin{minipage}[b]{\linewidth}\raggedright
Uso en el reporte
\end{minipage} \\
\midrule\noalign{}
\endhead
\bottomrule\noalign{}
\endlastfoot
1 & Black \& Wiliam, \emph{Assessment and Classroom Learning} & 1998 &
Evaluación formativa y aprendizaje &
\href{https://doi.org/10.1080/0969595980050102}{10.1080/0969595980050102}
& Fundamento de evaluación para el aprendizaje \\
2 & Nicol \& Macfarlane-Dick, \emph{Formative Assessment and
Self-Regulated Learning} & 2006 & Autorregulación y retroalimentación &
\href{https://doi.org/10.1080/03075070600572090}{10.1080/03075070600572090}
& Base para metacognición y agencia \\
3 & Hattie \& Timperley, \emph{The Power of Feedback} & 2007 &
Retroalimentación y logro &
\href{https://doi.org/10.3102/003465430298487}{10.3102/003465430298487}
& Diseño de feedback humano/IA \\
4 & Shute, \emph{Focus on Formative Feedback} & 2008 & Retroalimentación
formativa &
\href{https://doi.org/10.3102/0034654307313795}{10.3102/0034654307313795}
& Criterios de feedback eficaz \\
5 & Zawacki-Richter et al., \emph{Systematic review of research on
artificial intelligence applications in higher education} & 2019 &
Revisión sistemática de IA en educación superior &
\href{https://doi.org/10.1186/s41239-019-0171-0}{10.1186/s41239-019-0171-0}
& Mapa de aplicaciones: predicción, evaluación, tutoría,
personalización \\
6 & Crompton \& Burke, \emph{Artificial intelligence in higher
education: the state of the field} & 2023 & Estado del campo AIEd en
educación superior &
\href{https://doi.org/10.1186/s41239-023-00392-8}{10.1186/s41239-023-00392-8}
& Tendencias recientes y usos dominantes \\
7 & Kasneci et al., \emph{ChatGPT for good?} & 2023 & Oportunidades y
desafíos de LLM en educación &
\href{https://doi.org/10.1016/j.lindif.2023.102274}{10.1016/j.lindif.2023.102274}
& Riesgos y posibilidades pedagógicas de LLM \\
8 & Weber-Wulff et al., \emph{Testing of detection tools for
AI-generated text} & 2023 & Detectores de texto generado por IA &
\href{https://doi.org/10.1007/s40979-023-00146-z}{10.1007/s40979-023-00146-z}
& Crítica a la vigilancia automática \\
9 & Bahroun et al., \emph{Transforming Education} & 2023 & Revisión
bibliométrica de IA generativa en educación &
\href{https://doi.org/10.3390/su151712983}{10.3390/su151712983} &
Panorama de expansión de GenAI educativa \\
10 & Darvishi et al., \emph{Impact of AI assistance on student agency} &
2024 & Agencia estudiantil y asistencia de IA &
\href{https://doi.org/10.1016/j.compedu.2023.104967}{10.1016/j.compedu.2023.104967}
& Dependencia, autorregulación y feedback asistido \\
11 & Lee et al., \emph{The impact of generative AI on higher education
learning and teaching} & 2024 & Perspectivas docentes sobre GenAI &
\href{https://doi.org/10.1016/j.caeai.2024.100221}{10.1016/j.caeai.2024.100221}
& Ambigüedad institucional y rediseño docente \\
12 & Bond et al., \emph{A meta systematic review of artificial
intelligence in higher education} & 2024 & Metarrevisión de IA en
educación superior &
\href{https://doi.org/10.1186/s41239-023-00436-z}{10.1186/s41239-023-00436-z}
& Necesidad de ética, rigor y colaboración interdisciplinar \\
13 & Zhang et al., \emph{A Systematic Literature Review of Empirical
Research on Applying Generative AI in Education} & 2024 & Evidencia
empírica sobre GenAI en educación &
\href{https://doi.org/10.1007/s44366-024-0028-5}{10.1007/s44366-024-0028-5}
& Aprendizaje, feedback, metacognición y evaluación transformadora \\
14 & Bastani et al., \emph{Generative AI Can Harm Learning} & 2024 &
Experimento de campo sobre IA y aprendizaje matemático &
\href{https://doi.org/10.2139/ssrn.4895486}{10.2139/ssrn.4895486} &
Evidencia sobre rendimiento aparente y dependencia \\
15 & Perkins et al., \emph{The Artificial Intelligence Assessment Scale}
& 2024 & Marco de integración ética de IA en evaluación &
\href{https://doi.org/10.53761/q3azde36}{10.53761/q3azde36} & Niveles de
uso permitido de IA en evaluación \\
16 & Furze et al., \emph{The AI Assessment Scale in action} & 2024 &
Implementación piloto de AIAS &
\href{https://doi.org/10.14742/ajet.9434}{10.14742/ajet.9434} &
Evidencia institucional de rediseño evaluativo \\
17 & Lye \& Lim, \emph{Generative Artificial Intelligence in Tertiary
Education} & 2024 & Principios de rediseño evaluativo &
\href{https://doi.org/10.3390/educsci14060569}{10.3390/educsci14060569}
& Principios ``against, avoid, adopt'' y evaluación terciaria \\
18 & Zhao, Chapman \& Sabet, \emph{Generative AI and Educational
Assessments} & 2024 & Revisión sistemática de IA generativa y evaluación
&
\href{https://doi.org/10.70953/ERPv51.2412006}{10.70953/ERPv51.2412006}
& Síntesis empírica sobre evaluación con GenAI \\
19 & Kofinas, Tsay \& Pike, \emph{The impact of generative AI on
academic integrity of authentic assessments} & 2025 & Integridad
académica y evaluación auténtica &
\href{https://doi.org/10.1111/bjet.13585}{10.1111/bjet.13585} & Tensión
entre autenticidad, integridad y rediseño \\
\end{longtable}

\section{IA y aprendizaje: de la automatización a la mediación
cognitiva}\label{ia-y-aprendizaje-de-la-automatizaciuxf3n-a-la-mediaciuxf3n-cognitiva}

La literatura reciente confirma que la IA en educación no es un fenómeno
homogéneo. Zawacki-Richter et al.~(2019) organizan sus aplicaciones en
cuatro grandes áreas: perfilamiento y predicción, evaluación, sistemas
adaptativos/personalización y tutoría inteligente. Esta clasificación es
útil porque evita reducir la IA educativa a chatbots o generación de
texto. La IA ya interviene en la analítica del aprendizaje, la
recomendación de recursos, la retroalimentación automatizada, la
detección de patrones de riesgo y la personalización de trayectorias.

Crompton y Burke (2023) muestran que la investigación sobre IA en
educación superior se intensificó antes de ChatGPT, pero la irrupción de
los modelos generativos aceleró las preguntas pedagógicas. En su
revisión, los usos principales de IA se agrupan en evaluación,
predicción, asistentes de IA, tutoría inteligente y gestión del
aprendizaje. Esta taxonomía permite comprender que el problema no
consiste únicamente en impedir el fraude, sino en rediseñar las
condiciones bajo las cuales se aprende, se practica, se recibe feedback
y se demuestra competencia.

Desde la perspectiva del aprendizaje, la IA puede operar en tres
niveles:

\begin{enumerate}
\def\labelenumi{\arabic{enumi}.}
\tightlist
\item
  \textbf{Nivel instrumental:} apoya búsqueda, resumen, traducción,
  organización y generación preliminar de ideas.
\item
  \textbf{Nivel cognitivo:} ofrece andamiaje, pistas, ejemplos,
  preguntas socráticas, retroalimentación adaptativa y simulaciones.
\item
  \textbf{Nivel epistémico:} modifica qué se considera saber,
  argumentar, demostrar, programar, diseñar, resolver o crear en un
  dominio profesional.
\end{enumerate}

El riesgo aparece cuando el nivel instrumental reemplaza al nivel
cognitivo. En ese caso, el estudiante produce más rápido, pero piensa
menos; entrega un producto plausible, pero no consolida esquemas
internos de comprensión. Por tanto, el criterio institucional no debe
ser únicamente ``usar o no usar IA'', sino determinar si la IA opera
como \textbf{andamio temporal para aprender} o como \textbf{sustituto
estable del razonamiento}.

\section{Evidencia sobre rendimiento aparente y deuda
cognitiva}\label{evidencia-sobre-rendimiento-aparente-y-deuda-cognitiva}

El concepto de ``rendimiento aparente'' planteado en el documento fuente
encuentra respaldo empírico en Bastani et al.~(2024). En un experimento
de campo con cerca de mil estudiantes de secundaria en clases de
matemáticas, el acceso a GPT-4 mejoró el desempeño durante la práctica
asistida. Sin embargo, cuando la herramienta fue retirada, los
estudiantes que usaron una versión no guiada del sistema obtuvieron
peores resultados que quienes nunca tuvieron acceso a la IA. La
interpretación educativa es directa: la IA puede mejorar la ejecución
inmediata sin producir aprendizaje transferible.

Este resultado no justifica prohibiciones generales. El mismo estudio
muestra que una versión con salvaguardas pedagógicas redujo los efectos
negativos. La diferencia no radica solo en el modelo, sino en el
\textbf{diseño de la interacción}. Una IA que entrega respuestas finales
puede deteriorar el aprendizaje; una IA que ofrece pistas, exige
explicación, pregunta por el procedimiento y retrasa la solución puede
funcionar como tutor.

En términos pedagógicos, la deuda cognitiva aparece cuando:

\begin{itemize}
\tightlist
\item
  el estudiante delega la descomposición del problema;
\item
  no formula hipótesis propias antes de consultar la IA;
\item
  copia soluciones sin reconstruirlas;
\item
  no identifica errores o supuestos del modelo;
\item
  no puede resolver una tarea análoga sin asistencia;
\item
  confunde fluidez textual con comprensión disciplinar.
\end{itemize}

De aquí se deriva una regla de diseño: \textbf{todo uso de IA en
actividades de aprendizaje debe incluir un momento sin IA para comprobar
transferencia, autonomía y retención conceptual}.

\section{Agencia estudiantil, metacognición y
autorregulación}\label{agencia-estudiantil-metacogniciuxf3n-y-autorregulaciuxf3n}

La agencia estudiantil es una dimensión crítica. Darvishi et al.~(2024)
examinan si los estudiantes aprenden de la retroalimentación
personalizada generada por IA o si dependen de ella. Esta pregunta es
central porque la retroalimentación automatizada puede mejorar la
revisión de tareas, pero también desplazar procesos de autoevaluación si
el estudiante no participa activamente en la interpretación del
feedback.

La literatura clásica sobre evaluación formativa ya había establecido
que la retroalimentación eficaz no consiste en entregar correcciones,
sino en reducir la brecha entre desempeño actual y desempeño esperado.
Black y Wiliam (1998) muestran que la evaluación formativa puede generar
ganancias sustanciales de aprendizaje cuando se integra al proceso de
enseñanza. Nicol y Macfarlane-Dick (2006) agregan que el feedback debe
apoyar la autorregulación: clarificar metas, facilitar autoevaluación,
ofrecer información de calidad, fomentar diálogo, motivar, cerrar
brechas y proporcionar información útil al docente.

En ambientes con IA, estos principios adquieren mayor relevancia. La
retroalimentación generada por IA debe evaluarse mediante cuatro
criterios:

\begin{longtable}[]{@{}
  >{\raggedright\arraybackslash}p{(\linewidth - 4\tabcolsep) * \real{0.3333}}
  >{\raggedright\arraybackslash}p{(\linewidth - 4\tabcolsep) * \real{0.3333}}
  >{\raggedright\arraybackslash}p{(\linewidth - 4\tabcolsep) * \real{0.3333}}@{}}
\toprule\noalign{}
\begin{minipage}[b]{\linewidth}\raggedright
Criterio
\end{minipage} & \begin{minipage}[b]{\linewidth}\raggedright
Pregunta orientadora
\end{minipage} & \begin{minipage}[b]{\linewidth}\raggedright
Riesgo si se omite
\end{minipage} \\
\midrule\noalign{}
\endhead
\bottomrule\noalign{}
\endlastfoot
Exactitud disciplinar & ¿El feedback es técnicamente correcto? &
Aprendizaje de errores plausibles \\
Acción pedagógica & ¿Indica cómo mejorar, no solo qué está mal? &
Corrección superficial \\
Metacognición & ¿Obliga al estudiante a explicar su decisión? &
Dependencia pasiva \\
Transferencia & ¿Prepara para resolver tareas nuevas? & Rendimiento
localizado sin aprendizaje \\
\end{longtable}

La IA debe ser configurada como mediadora de autorregulación, no como
juez final ni como productor invisible de la respuesta.

\section{Evaluación del aprendizaje en la era de
IA}\label{evaluaciuxf3n-del-aprendizaje-en-la-era-de-ia}

\subsection{Insuficiencia del producto
final}\label{insuficiencia-del-producto-final}

El documento fuente afirma que el ensayo estático o el producto final
aislado pierden fuerza evaluativa. Esta afirmación es consistente con la
literatura reciente sobre IA generativa y evaluación. Cuando una tarea
se puede resolver con una instrucción genérica a un modelo, la evidencia
de aprendizaje se vuelve débil. No porque el estudiante necesariamente
haya actuado con deshonestidad, sino porque el instrumento ya no
distingue entre desempeño humano, asistencia legítima, delegación
parcial o automatización total.

Por tanto, la evaluación debe pasar de una lógica de
\textbf{verificación de autoría} a una lógica de \textbf{validación de
competencia}. La pregunta deja de ser ``¿quién escribió este texto?'' y
pasa a ser:

\begin{itemize}
\tightlist
\item
  ¿Puede el estudiante explicar sus decisiones?
\item
  ¿Puede justificar la elección de fuentes, datos, métodos o criterios?
\item
  ¿Puede detectar errores generados por IA?
\item
  ¿Puede transferir el procedimiento a un problema nuevo?
\item
  ¿Puede actuar competentemente en vivo, bajo restricciones realistas?
\end{itemize}

\subsection{Evaluación auténtica y vulnerabilidad a la
IA}\label{evaluaciuxf3n-autuxe9ntica-y-vulnerabilidad-a-la-ia}

La evaluación auténtica no queda automáticamente protegida frente a la
IA. Kofinas, Tsay y Pike (2025) muestran que la IA generativa también
afecta evaluaciones que tradicionalmente se consideraban más
resistentes, porque puede ayudar a producir análisis, planes, informes,
simulaciones textuales y productos multimodales. La autenticidad, por sí
sola, no basta; debe combinarse con trazabilidad, defensa, especificidad
contextual y evidencia de proceso.

Una evaluación auténtica resiliente requiere al menos cinco condiciones:

\begin{enumerate}
\def\labelenumi{\arabic{enumi}.}
\tightlist
\item
  \textbf{Contexto local o situado:} uso de datos, casos, restricciones
  o experiencias que no puedan ser resueltas por respuestas genéricas.
\item
  \textbf{Proceso visible:} registro de decisiones, descartes,
  versiones, prompts, fuentes y cambios.
\item
  \textbf{Juicio disciplinar:} identificación de errores, sesgos,
  supuestos y límites de la IA.
\item
  \textbf{Interacción evaluativa:} defensa oral, entrevista,
  demostración en vivo o revisión crítica.
\item
  \textbf{Transferencia:} tarea análoga, variación del caso o aplicación
  sin IA.
\end{enumerate}

\section{Limitaciones de la detección
automática}\label{limitaciones-de-la-detecciuxf3n-automuxe1tica}

La respuesta institucional centrada en detectores de texto generado por
IA es técnicamente débil. Weber-Wulff et al.~(2023) evaluaron
herramientas de detección y concluyeron que no son suficientemente
fiables para tomar decisiones punitivas de alto impacto. Además, la
paráfrasis, la traducción automática y la edición humana pueden alterar
significativamente el resultado de los detectores.

Esto implica que los detectores no deben ser usados como prueba única de
fraude académico. Su uso, si existe, debe limitarse a señal preliminar
de revisión pedagógica o administrativa, nunca como evidencia
concluyente. Una política institucional responsable debería privilegiar:

\begin{itemize}
\tightlist
\item
  declaración transparente del uso de IA;
\item
  rediseño de tareas;
\item
  entrevistas de validación;
\item
  comparación con evidencia longitudinal del estudiante;
\item
  rúbricas que evalúen proceso, fuentes y juicio;
\item
  procedimientos de debido proceso cuando exista sospecha.
\end{itemize}

\section{Rediseño de la evaluación: de la prohibición al contrato
pedagógico}\label{rediseuxf1o-de-la-evaluaciuxf3n-de-la-prohibiciuxf3n-al-contrato-pedaguxf3gico}

La literatura reciente converge en una conclusión: prohibir
genéricamente la IA es poco sostenible. Perkins et al.~(2024) proponen
la AI Assessment Scale como un marco para definir niveles de uso
permitido de IA según los resultados de aprendizaje esperados. Furze et
al.~(2024) reportan una implementación piloto en la que una escala
explícita ayudó a reducir incertidumbre, orientar prácticas docentes y
promover mayor transparencia.

La utilidad de este enfoque es que convierte el uso de IA en un
\textbf{contrato pedagógico}. Cada actividad debe aclarar:

\begin{itemize}
\tightlist
\item
  qué usos de IA están permitidos;
\item
  qué usos deben declararse;
\item
  qué partes deben ser realizadas sin IA;
\item
  qué evidencia de proceso debe entregarse;
\item
  cómo se evaluará la contribución humana;
\item
  qué consecuencias tiene la delegación no declarada.
\end{itemize}

Una escala institucional mínima podría adoptar cinco niveles:

\begin{longtable}[]{@{}
  >{\raggedleft\arraybackslash}p{(\linewidth - 6\tabcolsep) * \real{0.3077}}
  >{\raggedright\arraybackslash}p{(\linewidth - 6\tabcolsep) * \real{0.2308}}
  >{\raggedright\arraybackslash}p{(\linewidth - 6\tabcolsep) * \real{0.2308}}
  >{\raggedright\arraybackslash}p{(\linewidth - 6\tabcolsep) * \real{0.2308}}@{}}
\toprule\noalign{}
\begin{minipage}[b]{\linewidth}\raggedleft
Nivel
\end{minipage} & \begin{minipage}[b]{\linewidth}\raggedright
Condición de uso de IA
\end{minipage} & \begin{minipage}[b]{\linewidth}\raggedright
Evidencia exigida
\end{minipage} & \begin{minipage}[b]{\linewidth}\raggedright
Tipo de tarea apropiada
\end{minipage} \\
\midrule\noalign{}
\endhead
\bottomrule\noalign{}
\endlastfoot
0 & Sin IA & Producción en vivo o bajo supervisión & Exámenes de base,
pruebas diagnósticas, habilidades críticas mínimas \\
1 & IA para exploración & Registro breve de consulta & Lluvia de ideas,
búsqueda inicial, delimitación temática \\
2 & IA para retroalimentación & Versión inicial, feedback recibido y
mejora justificada & Ensayos, informes, código, diseño de propuestas \\
3 & IA como copiloto & Bitácora completa, prompts, decisiones y defensa
& Proyectos aplicados, análisis de datos, simulaciones \\
4 & IA como componente explícito del producto & Auditoría de modelo,
validación humana y reflexión crítica & Prototipos, agentes,
herramientas educativas, productos profesionales \\
\end{longtable}

Esta escala no debe aplicarse de forma uniforme. Debe alinearse con el
resultado de aprendizaje. Si el resultado es escritura académica
autónoma, el nivel permitido será bajo. Si el resultado es diseñar un
flujo de trabajo profesional asistido por IA, el nivel permitido será
alto, pero la evaluación deberá centrarse en auditoría, juicio y
validación.

\section{Índice de resiliencia evaluativa frente a
IA}\label{uxedndice-de-resiliencia-evaluativa-frente-a-ia}

Para operacionalizar el rediseño, se propone un índice simple de
resiliencia evaluativa frente a IA. No pretende ser una métrica
psicométrica definitiva, sino una herramienta diagnóstica para revisar
instrumentos.

{[} IRE\_\{IA\}=\frac{P+D+T+M+C}{A+1} {]}

Donde:

\begin{itemize}
\tightlist
\item
  (P): evidencia de proceso, en escala 0--3.
\item
  (D): defensa oral, entrevista o demostración en vivo, en escala 0--3.
\item
  (T): transferencia a una situación nueva, en escala 0--3.
\item
  (M): exigencia metacognitiva o reflexión sobre decisiones, en escala
  0--3.
\item
  (C): contextualización local, situada o basada en datos propios, en
  escala 0--3.
\item
  (A): facilidad de automatización por IA, en escala 0--3.
\end{itemize}

Interpretación sugerida:

\begin{longtable}[]{@{}
  >{\raggedleft\arraybackslash}p{(\linewidth - 4\tabcolsep) * \real{0.4000}}
  >{\raggedright\arraybackslash}p{(\linewidth - 4\tabcolsep) * \real{0.3000}}
  >{\raggedright\arraybackslash}p{(\linewidth - 4\tabcolsep) * \real{0.3000}}@{}}
\toprule\noalign{}
\begin{minipage}[b]{\linewidth}\raggedleft
Valor de (IRE\_\{IA\})
\end{minipage} & \begin{minipage}[b]{\linewidth}\raggedright
Lectura
\end{minipage} & \begin{minipage}[b]{\linewidth}\raggedright
Acción recomendada
\end{minipage} \\
\midrule\noalign{}
\endhead
\bottomrule\noalign{}
\endlastfoot
0--2 & Alta vulnerabilidad & Rediseñar antes de aplicar \\
\textgreater2--5 & Vulnerabilidad media & Añadir trazabilidad, defensa o
transferencia \\
\textgreater5--10 & Resiliencia aceptable & Aplicar con rúbrica
explícita \\
\textgreater10 & Alta resiliencia & Usar en evaluación de alto impacto,
con control de carga docente \\
\end{longtable}

Este índice puede utilizarse en comités curriculares para auditar tareas
existentes. Por ejemplo, un ensayo genérico sin defensa puede tener alto
(A) y bajo (P, D, T, M, C). En cambio, un portafolio con bitácora, datos
locales, defensa oral y tarea de transferencia reduce el riesgo de
delegación invisible.

\section{Taxonomía post-IA del aprendizaje
evaluable}\label{taxonomuxeda-post-ia-del-aprendizaje-evaluable}

El documento fuente propone desplazar la taxonomía desde verbos de bajo
nivel hacia acciones como interpretar, contextualizar, integrar,
adaptar, criticar, interrogar, deliberar e innovar. Esta idea puede
expandirse en una taxonomía de desempeño evaluable:

\begin{longtable}[]{@{}
  >{\raggedright\arraybackslash}p{(\linewidth - 4\tabcolsep) * \real{0.3333}}
  >{\raggedright\arraybackslash}p{(\linewidth - 4\tabcolsep) * \real{0.3333}}
  >{\raggedright\arraybackslash}p{(\linewidth - 4\tabcolsep) * \real{0.3333}}@{}}
\toprule\noalign{}
\begin{minipage}[b]{\linewidth}\raggedright
Dominio
\end{minipage} & \begin{minipage}[b]{\linewidth}\raggedright
Pregunta evaluativa
\end{minipage} & \begin{minipage}[b]{\linewidth}\raggedright
Evidencia robusta
\end{minipage} \\
\midrule\noalign{}
\endhead
\bottomrule\noalign{}
\endlastfoot
Comprensión & ¿El estudiante puede explicar el concepto sin reproducir
texto generado? & Explicación oral, mapa conceptual, ejemplo propio \\
Aplicación & ¿Puede usar el concepto en un caso nuevo? & Resolución
contextualizada, práctica supervisada \\
Crítica & ¿Puede detectar errores, sesgos o supuestos? & Auditoría de
respuesta IA, dictamen técnico \\
Integración & ¿Puede combinar fuentes, datos y métodos? & Portafolio,
informe con trazabilidad \\
Creación & ¿Puede producir una solución justificada? & Prototipo,
diseño, simulación, producto aplicado \\
Metacognición & ¿Puede explicar cómo aprendió y qué cambió? & Bitácora
reflexiva, defensa del proceso \\
Autonomía & ¿Puede desempeñarse sin IA cuando es necesario? & Tarea
espejo sin asistencia, entrevista técnica \\
\end{longtable}

La clave no es eliminar la IA, sino preservar espacios donde la
competencia humana sea observable.

\section{Propuesta de modelo institucional de evaluación
resiliente}\label{propuesta-de-modelo-institucional-de-evaluaciuxf3n-resiliente}

Se propone un modelo en cinco fases.

\subsection{Fase 1. Auditoría de
vulnerabilidad}\label{fase-1.-auditoruxeda-de-vulnerabilidad}

Cada instrumento debe ser sometido a una prueba básica: intentar
resolverlo con una IA generativa común. Si la IA produce una respuesta
de alta calidad con una instrucción genérica, el instrumento es
vulnerable. No debe aplicarse sin ajustes.

\subsection{Fase 2. Definición del nivel de IA
permitido}\label{fase-2.-definiciuxf3n-del-nivel-de-ia-permitido}

El docente debe declarar el nivel de IA permitido antes de la actividad.
La regla debe estar en la guía, la rúbrica y las instrucciones de
entrega. La ambigüedad aumenta el conflicto ético y reduce la validez de
la evaluación.

\subsection{Fase 3. Diseño de evidencia
múltiple}\label{fase-3.-diseuxf1o-de-evidencia-muxfaltiple}

La calificación no debe depender de una sola evidencia textual. Una
evaluación resiliente combina:

\begin{itemize}
\tightlist
\item
  producto final;
\item
  bitácora de proceso;
\item
  fuentes verificadas;
\item
  explicación de decisiones;
\item
  revisión de errores;
\item
  defensa oral o entrevista;
\item
  transferencia a caso nuevo.
\end{itemize}

\subsection{Fase 4. Rúbrica orientada al juicio
humano}\label{fase-4.-ruxfabrica-orientada-al-juicio-humano}

La rúbrica debe separar el valor del producto y el valor del proceso. Un
producto correcto con proceso opaco no debe recibir la máxima
calificación en actividades de alto impacto. La rúbrica debe incluir
criterios como trazabilidad, justificación, crítica de IA,
contextualización y autonomía.

\subsection{Fase 5. Retroalimentación para aprendizaje, no solo
calificación}\label{fase-5.-retroalimentaciuxf3n-para-aprendizaje-no-solo-calificaciuxf3n}

La evaluación debe cerrar el ciclo formativo. La IA puede apoyar
feedback rápido, pero el docente conserva la responsabilidad de validar
criterios, ajustar orientaciones y atender errores conceptuales. La
retroalimentación debe promover revisión, no simple corrección
cosmética.

\section{Implicaciones para docentes}\label{implicaciones-para-docentes}

El rol docente se desplaza hacia diseño, curaduría, validación y
acompañamiento metacognitivo. Las competencias docentes prioritarias
son:

\begin{itemize}
\tightlist
\item
  diseñar tareas con distintos niveles de uso permitido de IA;
\item
  formular rúbricas que valoren proceso y transferencia;
\item
  construir defensas orales breves y sostenibles;
\item
  evaluar bitácoras de IA sin convertirlas en carga administrativa
  excesiva;
\item
  detectar errores plausibles en respuestas generadas;
\item
  enseñar a contrastar fuentes y validar afirmaciones;
\item
  diferenciar entre asistencia legítima y delegación indebida;
\item
  proteger datos personales al usar plataformas externas.
\end{itemize}

El docente no necesita convertirse en ingeniero de IA, pero sí debe
comprender sus efectos sobre el aprendizaje, la autoría, el feedback y
la validez evaluativa.

\section{Implicaciones para
estudiantes}\label{implicaciones-para-estudiantes}

El estudiante debe aprender a usar IA sin perder agencia. Esto exige
normas explícitas y entrenamiento progresivo. Las competencias mínimas
son:

\begin{itemize}
\tightlist
\item
  declarar el uso de IA con honestidad;
\item
  formular prompts orientados al aprendizaje, no a la sustitución;
\item
  pedir pistas antes que respuestas completas;
\item
  contrastar salidas con fuentes académicas;
\item
  detectar sesgos, errores y alucinaciones;
\item
  explicar qué parte del producto es propia y qué parte fue asistida;
\item
  demostrar transferencia sin IA cuando el resultado de aprendizaje lo
  requiera.
\end{itemize}

Una alfabetización crítica en IA no se limita a saber usar herramientas.
Incluye saber cuándo no usarlas.

\section{Recomendaciones operativas}\label{recomendaciones-operativas}

\begin{enumerate}
\def\labelenumi{\arabic{enumi}.}
\tightlist
\item
  \textbf{Crear una política institucional de uso de IA por niveles.}
  Cada actividad debe indicar explícitamente el nivel permitido de IA.
\item
  \textbf{Eliminar evaluaciones genéricas de bajo contexto.} Preguntas
  puramente definicionales o ensayos universales son altamente
  automatizables.
\item
  \textbf{Introducir bitácoras breves de IA.} No deben ser documentos
  extensos; basta con registrar herramienta, propósito, prompts
  principales, decisiones y validaciones.
\item
  \textbf{Incluir defensas orales proporcionales.} No toda actividad
  requiere defensa completa, pero las evaluaciones de alto impacto
  deberían incluir al menos una entrevista breve o explicación en vivo.
\item
  \textbf{Usar tareas espejo.} Una parte del proceso puede realizarse
  con IA y otra sin IA, para verificar transferencia.
\item
  \textbf{No usar detectores como evidencia concluyente.} Pueden
  orientar revisión, pero no reemplazan debido proceso ni juicio
  docente.
\item
  \textbf{Capacitar en feedback asistido por IA.} El estudiante debe
  aprender a solicitar retroalimentación útil y a evaluarla
  críticamente.
\item
  \textbf{Proteger datos personales.} Las guías deben prohibir
  introducir información sensible de estudiantes, docentes, pacientes,
  usuarios o instituciones en sistemas externos sin autorización.
\item
  \textbf{Alinear evaluación con resultados de aprendizaje.} La pregunta
  central es qué competencia humana se certifica y qué tipo de IA es
  compatible con dicha certificación.
\item
  \textbf{Revisar periódicamente los instrumentos.} La resiliencia
  evaluativa no es permanente; cambia con la capacidad de los modelos.
\end{enumerate}

\section{Conclusiones}\label{conclusiones}

La IA generativa no elimina la evaluación, pero sí invalida buena parte
de las prácticas evaluativas centradas en productos textuales
descontextualizados. El problema no es solo la autoría; es la validez de
la inferencia: a partir de una entrega, ¿podemos concluir que el
estudiante aprendió? En muchos casos, la respuesta ya no es afirmativa
si el instrumento no exige proceso, defensa, transferencia y
metacognición.

La literatura revisada permite sostener una posición equilibrada. La IA
puede ampliar oportunidades de aprendizaje mediante tutoría,
retroalimentación, personalización y apoyo cognitivo. Sin embargo,
también puede producir dependencia, rendimiento aparente, pérdida de
agencia y debilitamiento de habilidades si se usa como sustituto del
esfuerzo intelectual. Por tanto, el desafío central es diseñar
interacciones y evaluaciones donde la IA funcione como andamio, no como
reemplazo.

La institución educativa debe pasar de la vigilancia reactiva al diseño
pedagógico proactivo. Una evaluación resiliente frente a IA no se basa
en perseguir estudiantes, sino en construir evidencias robustas de
competencia humana: explicar, justificar, contrastar, transferir,
deliberar y actuar con criterio. La autonomía humana no se protege
prohibiendo toda tecnología; se protege diseñando condiciones en las que
el estudiante deba demostrar que puede pensar con IA, contra la IA y sin
IA cuando sea necesario.

\section{Referencias}\label{referencias}

Black, P., \& Wiliam, D. (1998). Assessment and classroom learning.
\emph{Assessment in Education: Principles, Policy \& Practice, 5}(1),
7--74. https://doi.org/10.1080/0969595980050102

Nicol, D. J., \& Macfarlane-Dick, D. (2006). Formative assessment and
self-regulated learning: A model and seven principles of good feedback
practice. \emph{Studies in Higher Education, 31}(2), 199--218.
https://doi.org/10.1080/03075070600572090

Hattie, J., \& Timperley, H. (2007). The power of feedback. \emph{Review
of Educational Research, 77}(1), 81--112.
https://doi.org/10.3102/003465430298487

Shute, V. J. (2008). Focus on formative feedback. \emph{Review of
Educational Research, 78}(1), 153--189.
https://doi.org/10.3102/0034654307313795

Zawacki-Richter, O., Marín, V. I., Bond, M., \& Gouverneur, F. (2019).
Systematic review of research on artificial intelligence applications in
higher education: Where are the educators? \emph{International Journal
of Educational Technology in Higher Education, 16}, Article 39.
https://doi.org/10.1186/s41239-019-0171-0

Crompton, H., \& Burke, D. (2023). Artificial intelligence in higher
education: The state of the field. \emph{International Journal of
Educational Technology in Higher Education, 20}, Article 22.
https://doi.org/10.1186/s41239-023-00392-8

Kasneci, E., Sessler, K., Küchemann, S., Bannert, M., Dementieva, D.,
Fischer, F., et al.~(2023). ChatGPT for good? On opportunities and
challenges of large language models for education. \emph{Learning and
Individual Differences, 103}, 102274.
https://doi.org/10.1016/j.lindif.2023.102274

Weber-Wulff, D., Anohina-Naumeca, A., Bjelobaba, S., Foltýnek, T.,
Guerrero-Dib, J., Popoola, O., Šigut, P., \& Waddington, L. (2023).
Testing of detection tools for AI-generated text. \emph{International
Journal for Educational Integrity, 19}, Article 26.
https://doi.org/10.1007/s40979-023-00146-z

Bahroun, Z., Anane, C., Ahmed, V., \& Zacca, A. (2023). Transforming
education: A comprehensive review of generative artificial intelligence
in educational settings through bibliometric and content analysis.
\emph{Sustainability, 15}(17), 12983.
https://doi.org/10.3390/su151712983

Darvishi, A., Khosravi, H., Sadiq, S., Gašević, D., \& Siemens, G.
(2024). Impact of AI assistance on student agency. \emph{Computers \&
Education, 210}, 104967. https://doi.org/10.1016/j.compedu.2023.104967

Lee, D., Arnold, M., Srivastava, A., Plastow, K., Strelan, P., Ploeckl,
F., Lekkas, D., \& Palmer, E. (2024). The impact of generative AI on
higher education learning and teaching: A study of educators'
perspectives. \emph{Computers and Education: Artificial Intelligence,
6}, 100221. https://doi.org/10.1016/j.caeai.2024.100221

Bond, M., Khosravi, H., De Laat, M., Bergdahl, N., Negrea, V., Oxley,
E., Pham, P., Chong, S. W., \& Siemens, G. (2024). A meta systematic
review of artificial intelligence in higher education: A call for
increased ethics, collaboration, and rigour. \emph{International Journal
of Educational Technology in Higher Education, 21}, Article 4.
https://doi.org/10.1186/s41239-023-00436-z

Zhang, X., Zhang, P., Shen, Y., Liu, M., Wang, Q., Gašević, D., \& Fan,
Y. (2024). A systematic literature review of empirical research on
applying generative artificial intelligence in education.
\emph{Frontiers of Digital Education, 1}(3), 223--245.
https://doi.org/10.1007/s44366-024-0028-5

Bastani, H., Bastani, O., Sungu, A., Ge, H., Kabakcı, Ö., \& Mariman, R.
(2024). \emph{Generative AI can harm learning}. The Wharton School
Research Paper. https://doi.org/10.2139/ssrn.4895486

Perkins, M., Furze, L., Roe, J., \& MacVaugh, J. (2024). The Artificial
Intelligence Assessment Scale (AIAS): A framework for ethical
integration of generative AI in educational assessment. \emph{Journal of
University Teaching and Learning Practice, 21}(06).
https://doi.org/10.53761/q3azde36

Furze, L., Perkins, M., Roe, J., \& MacVaugh, J. (2024). The AI
Assessment Scale (AIAS) in action: A pilot implementation of
GenAI-supported assessment. \emph{Australasian Journal of Educational
Technology, 40}(4), 38--55. https://doi.org/10.14742/ajet.9434

Lye, C. Y., \& Lim, L. (2024). Generative artificial intelligence in
tertiary education: Assessment redesign principles and considerations.
\emph{Education Sciences, 14}(6), 569.
https://doi.org/10.3390/educsci14060569

Zhao, J., Chapman, E., \& Ghassemi Pour Sabet, P. (2024). Generative AI
and educational assessments: A systematic review. \emph{Education
Research \& Perspectives, 51}, 124--155.
https://doi.org/10.70953/ERPv51.2412006

Kofinas, A. K., Tsay, C. H.-H., \& Pike, D. (2025). The impact of
generative AI on academic integrity of authentic assessments within a
higher education context. \emph{British Journal of Educational
Technology, 56}(6), 2522--2549. https://doi.org/10.1111/bjet.13585




\end{document}
